\section{Method}
\label{sec:method}

\textsc{VisInject} decouples \emph{generating} the attack signal from \emph{transporting} it onto a natural photo, and from \emph{measuring} the result. Figure~\ref{fig:pipeline} shows the three stages.

\begin{figure}[t]
  \centering
  \resizebox{\textwidth}{!}{% pipeline.tex — TikZ block diagram of the three-stage VisInject attack pipeline.

\begin{tikzpicture}[
  font=\sffamily\small,
  >=Latex,
  every node/.style={align=center},
  stage/.style={
    rectangle, rounded corners=4pt,
    draw=NavyAccent, line width=0.6pt,
    fill=CardBG,
    minimum width=4.5cm, minimum height=2.4cm,
    inner sep=6pt,
  },
  arrow/.style={->, line width=1pt, draw=NavyAccent},
  caption/.style={font=\sffamily\footnotesize\itshape, color=SubInk},
]

% Three stages
\node[stage] (s1) {%
  \textbf{Stage 1}\\
  \textbf{\color{NavyAccent} UniversalAttack}\\[2pt]
  \footnotesize PGD on grey image\\
  \footnotesize 60 benign prompts\\
  \footnotesize Multi-VLM loss sum
};
\node[stage, right=2.0cm of s1] (s2) {%
  \textbf{Stage 2}\\
  \textbf{\color{AccentBlue} AnyAttack Fusion}\\[2pt]
  \footnotesize CLIP encodes universal img\\
  \footnotesize Decoder $\to$ \eps-bounded noise\\
  \footnotesize Add to any clean photo
};
\node[stage, right=2.0cm of s2] (s3) {%
  \textbf{Stage 3}\\
  \textbf{\color{GoodGreen} Dual-dim Evaluation}\\[2pt]
  \footnotesize Output Affected check\\
  \footnotesize Target Injected check\\
  \footnotesize Programmatic, no API
};

% Arrows
\draw[arrow] (s1) -- (s2)
  node[midway, above, caption] {universal\\image}
  node[midway, below, caption] {(once per\\target phrase)};
\draw[arrow] (s2) -- (s3)
  node[midway, above, caption] {adversarial\\photo}
  node[midway, below, caption] {(once per clean\\image)};

% Inputs row
\node[below=0.65cm of s1, caption] (in1) {\textbf{Input:} target phrase\\\& white-box VLMs};
\draw[arrow] (in1) -- (s1);
\node[below=0.65cm of s2, caption] (in2) {\textbf{Input:} clean photo};
\draw[arrow] (in2) -- (s2);
\node[below=0.65cm of s3, caption] (in3) {\textbf{Input:} target VLM\\\& 45 benign questions};
\draw[arrow] (in3) -- (s3);

% Output
\node[right=1.2cm of s3, caption, text=BadRed] (out) {Disruption rate\\Injection rate};
\draw[arrow] (s3) -- (out);

\end{tikzpicture}
}
  \caption{The three-stage \textsc{VisInject} pipeline. Stage~1 trains a single universal adversarial image against $N$ white-box \vlm{}s; Stage~2 transports the resulting signal onto an arbitrary clean photo using a pretrained \textsc{AnyAttack} decoder under an \Linf{} budget; Stage~3 evaluates each (clean, adversarial) response pair along two independent axes.}
  \label{fig:pipeline}
\end{figure}

\subsection{Stage 1 --- UniversalAttack}
\label{sec:method-stage1}

Let $f_i: \mathcal{I} \times \mathcal{T} \to \mathcal{T}$ denote the $i$-th surrogate \vlm{}, $x_0 \in [0,1]^{H \times W \times 3}$ a fixed grey image, $\mathcal{P} = \{p_1, \ldots, p_{60}\}$ a set of $60$ benign questions, and $y^\star$ the chosen target phrase. We optimise an unbounded latent $z_1 \in \mathbb{R}^{H \times W \times 3}$ with the reparameterisation
\begin{equation}
  x = 0.5 + \gamma \cdot \tanh(z_1), \qquad \gamma = 0.5,
  \label{eq:reparam}
\end{equation}
which keeps $x \in [0,1]$ without an explicit projection step. The training objective is the sum of token-level cross-entropy losses across all surrogates and all benign prompts:
\begin{equation}
  \mathcal{L}(x) \;=\; \sum_{i=1}^{N} \sum_{p \in \mathcal{P}} \mathrm{CE}\!\bigl(f_i(x, p),\; y^\star\bigr).
  \label{eq:loss}
\end{equation}
We minimise $\mathcal{L}$ with Adam ($\text{lr}=10^{-2}$) for $2{,}000$ steps. The output is a single \emph{universal adversarial image} $x_u$ that nudges every white-box \vlm{} toward $y^\star$ regardless of the question. The surrogate ensemble size $N \in \{2, 3, 4\}$ is the ``model configuration'' (\texttt{2m}, \texttt{3m}, \texttt{4m}) that we sweep in \S\ref{sec:experiments}.

\subsection{Stage 2 --- AnyAttack Fusion}
\label{sec:method-stage2}

The universal image $x_u$ is itself adversarial, but it does not look like a real photo, so it would betray the attack to a human looking at the upload. Stage~2 transports the attack signal onto an arbitrary clean photo $x_c$ using the pretrained \textsc{AnyAttack} decoder of \citet{zhang2025anyattack}, weights \texttt{coco\_bi.pt}.

A frozen CLIP ViT-B/16 \citep{radford2021clip} encodes $x_u$ into a 768-dimensional feature, which the \textsc{AnyAttack} decoder maps to a noise tensor $\delta(x_u)$ of the same spatial size as $x_c$. We add this noise under an \Linf{} budget:
\begin{equation}
  x_a \;=\; \mathrm{clip}_{[0,1]}\!\Bigl(x_c + \mathrm{clip}_{[-\eps,\eps]}\bigl(\delta(x_u)\bigr)\Bigr), \qquad \eps = 16/255.
  \label{eq:fuse}
\end{equation}
Crucially, the decoder is reused as-is~--~we do not retrain it. This makes the attack reproducible from public weights and decouples Stage~1 (which is expensive) from Stage~2 (which is essentially free per new clean image). The cost is that the decoder sometimes erases payload specifics --- a tradeoff we quantify in \S\ref{sec:results}.

\subsection{Stage 3 --- Dual-Dimension Evaluation}
\label{sec:method-stage3}

Given a clean image $x_c$ and its adversarial counterpart $x_a$, we feed both to a target \vlm{} together with each of $45$ benign questions and collect the answers $r_c$ and $r_a$. A first-version evaluation \citep{rahmatullaev2025universal} would ask a judge model to rate ``how much $r_a$ deviates from $r_c$''. We found this conflates two effects:

\begin{itemize}
  \item \textbf{Output Affected (drift).} Did $r_a$ differ in any meaningful way from $r_c$? This catches the \vlm{} producing nonsense, hallucinated text fragments, or off-topic descriptions --- common, easy to score, but \emph{not} the same as a successful attack.
  \item \textbf{Target Injected (payload).} Did $r_a$ contain the actual target phrase $y^\star$ (or a recognisable variant)? This is the metric the threat model in \S\ref{sec:threat} actually cares about.
\end{itemize}

Our second-version evaluation runs two independent programmatic checks on every (clean, adversarial) response pair: a longest-common-subsequence drift score for Output Affected, and a keyword/regex match for Target Injected. The dual-axis judgements remove the ``a deviation $\Rightarrow$ an injection'' shortcut and give us the gap reported in \S\ref{sec:results}.
